\documentclass[11pt]{article}

% ---------------------------------------------------------------------------
% AI Applications in Research, with a Focus on Claude Code
% A concept summary compiled for Scott's Mixtape readers.
% Compiles with pdflatex or xelatex using standard TeX Live packages.
% ---------------------------------------------------------------------------

\usepackage[utf8]{inputenc}
\usepackage[T1]{fontenc}
\usepackage[margin=1in]{geometry}
\usepackage{microtype}
\usepackage{enumitem}
\usepackage{booktabs}
\usepackage{xcolor}
\usepackage{listings}
\usepackage[colorlinks=true,linkcolor=black,urlcolor=blue,citecolor=black]{hyperref}

\definecolor{codebg}{gray}{0.96}
\lstset{
  basicstyle=\ttfamily\small,
  backgroundcolor=\color{codebg},
  frame=single,
  rulecolor=\color{gray!40},
  breaklines=true,
  columns=fullflexible,
  keepspaces=true,
  showstringspaces=false,
  xleftmargin=0.5em,
  xrightmargin=0.5em,
}

\title{\textbf{AI Applications in Research}\\[2pt]
\large A Concept Summary, with a Focus on Agentic Coding and Claude Code}
\author{Compiled for \emph{Scott's Mixtape Substack} readers}
\date{September 21, 2026}

\begin{document}
\maketitle

\begin{abstract}
\noindent This document summarizes the main ways artificial intelligence---and
large language models (LLMs) in particular---are being used across the research
workflow, from ideation to replication. It gives special attention to
\emph{agentic coding assistants}, using Anthropic's \textbf{Claude Code} as the
running example, and to the methodological guardrails that keep AI-assisted
research valid and reproducible. A short case study revisits the
Ludwig--Mullainathan--Rambachan (2026) econometric framework for LLMs, which was
the subject of a recent \emph{Scott's Mixtape} explainer.
\end{abstract}

\tableofcontents
\bigskip

% ---------------------------------------------------------------------------
\section{Introduction}
% ---------------------------------------------------------------------------
LLMs have moved from a curiosity to a general-purpose instrument in empirical
research. Two shifts drive this. First, models can now read, write, and reason
over text and code at a scale and cost that make previously infeasible projects
routine. Second, \emph{agentic} tools---systems that can plan, call other tools,
edit files, and execute commands---let researchers delegate multi-step technical
work rather than only asking for isolated answers. This document maps where AI
adds value in the research pipeline, then focuses on the agentic coding case
(Claude Code) and the conditions under which AI-assisted results remain
scientifically valid.

% ---------------------------------------------------------------------------
\section{Where AI Enters the Research Workflow}
% ---------------------------------------------------------------------------
AI applications span the full arc of a project. The table below organizes them by
stage, with the caveat that the highest-leverage uses are also the ones that most
require human verification.

\begin{table}[h]
\centering
\small
\begin{tabular}{@{}p{3.1cm}p{7.4cm}p{3.3cm}@{}}
\toprule
\textbf{Stage} & \textbf{Representative AI uses} & \textbf{Primary risk} \\
\midrule
Ideation \& scoping & Brainstorming hypotheses, mapping a literature, stress-testing identification strategies. & Plausible-but-wrong framing. \\
Literature review & Summarizing papers, extracting claims, building citation maps. & Hallucinated references. \\
Data collection & Scraping, parsing PDFs, structuring unstructured text. & Silent parsing errors. \\
Measurement / labeling & Coding text into economic variables (sentiment, topics, categories) at scale. & Non-classical measurement error. \\
Analysis \& coding & Writing, refactoring, and debugging analysis code; building pipelines. & Subtle logic bugs. \\
Writing & Drafting, editing, tightening prose; formatting \LaTeX{} and tables. & Overconfident claims. \\
Replication & Reconstructing pipelines, porting code across languages, writing tests. & Undetected divergence. \\
Review & Checking robustness, generating counterarguments, sanity-checking results. & False reassurance. \\
\bottomrule
\end{tabular}
\caption{AI uses across the research workflow, and the dominant failure mode at each stage.}
\end{table}

\paragraph{The measurement case is special.} When an LLM \emph{labels} data whose
labels then enter a regression, the tool stops being a convenience and becomes
part of the estimator. This is where formal econometric care is required
(Section~\ref{sec:casestudy}).

% ---------------------------------------------------------------------------
\section{Agentic Coding Assistants and Claude Code}
% ---------------------------------------------------------------------------
A chat assistant answers questions; an \emph{agentic} assistant does work. It
operates inside your project, reads and edits files, runs commands, inspects the
results, and iterates toward a goal you specify in natural language.

\subsection{What Claude Code Is}
\textbf{Claude Code} is Anthropic's command-line (and IDE / web) agent for
software and data work. Given a task, it can:
\begin{itemize}[leftmargin=1.4em,itemsep=2pt]
  \item explore an unfamiliar codebase and explain how it fits together;
  \item write, edit, and refactor code across many files;
  \item run tests, linters, and analysis scripts, then read the output and fix what broke;
  \item perform version-control operations (branches, commits, pull requests);
  \item extend its own reach through the Model Context Protocol (MCP) and specialized subagents.
\end{itemize}
The interaction model is a loop: \emph{you state intent $\rightarrow$ the agent
plans and acts $\rightarrow$ you inspect the diff and results $\rightarrow$ you
refine.} The human stays in the loop as reviewer and director.

\subsection{Why It Fits Research Computing}
Research code is often a thicket of scripts glued together under deadline
pressure. Agentic assistance helps most where that friction lives:
\begin{description}[leftmargin=1.6em,style=nextline,itemsep=3pt]
  \item[Reproducible pipelines] Turn ad-hoc scripts into an ordered, documented, re-runnable pipeline (e.g., a \texttt{Makefile} or a numbered script sequence) with a single clean entry point.
  \item[Code archaeology] Onboard onto a collaborator's or your own past project by asking the agent to trace data flow from raw inputs to final tables.
  \item[Language ports] Translate an estimator or simulation between R, Python, Stata, and Julia while preserving numerical behavior, then diff the outputs.
  \item[Testing] Generate unit tests and regression tests around analysis functions so that later edits cannot silently change results.
  \item[Writing \& \LaTeX{}] Draft, restructure, and format papers, tables, and figures; wire up captions, labels, and references.
  \item[Data wrangling] Parse messy PDFs, harmonize inconsistent files, and build cleaning steps that are inspectable rather than opaque.
\end{description}

\subsection{Working Practices That Pay Off}
\begin{enumerate}[leftmargin=1.6em,itemsep=3pt]
  \item \textbf{Write a project brief the agent can read.} A \texttt{CLAUDE.md} (or equivalent) file describing conventions, data locations, and ``never do this'' rules keeps behavior consistent across sessions.
  \item \textbf{Prefer small, verifiable steps.} Ask for one change at a time and inspect the diff; a validated small push beats a sprawling unreviewed one.
  \item \textbf{Make the agent prove its work.} Have it run the tests/linters and show passing output before you accept a change---and reproduce a bug before ``fixing'' it.
  \item \textbf{Pin the environment.} Lockfiles, fixed seeds, and recorded package versions are what make an AI-built pipeline reproducible by others.
  \item \textbf{Keep secrets out.} Do not paste credentials or private data into prompts or committed files.
  \item \textbf{Treat the agent as a fast junior collaborator, not an oracle.} It accelerates the mechanical work; you own the scientific judgment.
\end{enumerate}

\subsection{A Minimal Workflow Sketch}
A typical Claude Code loop for a research task reads like a conversation with a
capable assistant who then goes and does the work:

\begin{lstlisting}[language=bash]
# 1. Point the agent at the task, in plain language:
"Refactor analysis/clean.R into an idempotent pipeline,
 add tests for the winsorization step, and keep outputs
 byte-identical to the current results/ directory."

# 2. The agent explores, edits, and runs the checks itself:
Rscript analysis/clean.R
Rscript tests/test_clean.R        # <- agent reads pass/fail

# 3. You review the diff, then let it commit:
git add -A && git commit -m "Refactor cleaning into tested pipeline"
\end{lstlisting}

The point is not that the agent replaces the researcher, but that it collapses
the distance between intent (``make this reproducible and tested'') and a
concrete, inspectable diff.

% ---------------------------------------------------------------------------
\section{Methodological Guardrails}
% ---------------------------------------------------------------------------
AI accelerates production; it does not by itself guarantee validity. Four
recurring hazards deserve explicit management.

\begin{description}[leftmargin=1.6em,style=nextline,itemsep=3pt]
  \item[Hallucination] Generated text and citations can be confidently wrong; verify every factual claim and reference against primary sources.
  \item[Training leakage] If a model has effectively ``seen'' the answer (or the test data) during training, apparent predictive success is illusory; control it through model choice and design.
  \item[Non-classical measurement error] LLM labels can be wrong in ways correlated with the truth, biasing coefficients in any direction---including sign flips---so high accuracy alone is not a defense.
  \item[Reproducibility drift] Model versions, prompts, and sampling temperature all move results; record them as you would any other part of a research protocol.
\end{description}

% ---------------------------------------------------------------------------
\section{Case Study: The Econometrics of LLMs}
\label{sec:casestudy}
% ---------------------------------------------------------------------------
A recent \emph{Scott's Mixtape} post explains Ludwig, Mullainathan, and
Rambachan's \emph{``Large Language Models: An Applied Econometric Framework''}
(Annual Review of Economics, 2026). The framework is the cleanest current answer
to ``when may I use an LLM in an empirical study?'' and it divides the question in
two.

\paragraph{Prediction problems.} When the goal is to predict an outcome, validity
requires \textbf{no training leakage} between the LLM's training data and the
researcher's evaluation sample. This is enforceable through careful model
selection and research design.

\paragraph{Estimation (measurement) problems.} When an LLM \emph{measures} an
economic concept whose values then enter downstream analysis, valid inference
requires pairing LLM outputs with a \textbf{small human-coded validation sample}.
That validation sample is what makes it possible to estimate---and correct
for---the LLM's labeling error.

\paragraph{Why the validation sample is non-negotiable.} Absent it, researchers
cannot assess the error in LLM outputs, and seemingly innocuous choices (which
model, which prompt) can produce \emph{dramatically different} parameter
estimates. Because LLM measurement error is non-classical, raw accuracy does not
bound the bias; a modest, well-designed validation set does. The practical rule
for AI-assisted measurement: spend a little human labeling to debias a lot of
cheap machine labeling.

% ---------------------------------------------------------------------------
\section{Conclusion}
% ---------------------------------------------------------------------------
AI, and agentic tools such as Claude Code in particular, compress the mechanical
cost of research: exploring code, building reproducible pipelines, testing,
porting, drafting, and formatting. The gains are real and large. But the same
tools introduce failure modes---hallucination, leakage, non-classical
measurement error, and reproducibility drift---that require the researcher to
stay firmly in the loop as verifier and author. Used with a project brief, small
verifiable steps, pinned environments, and (for measurement) a validation
sample, AI becomes a dependable collaborator rather than a source of hidden risk.

% ---------------------------------------------------------------------------
\section*{References \& Resources}
% ---------------------------------------------------------------------------
\addcontentsline{toc}{section}{References \& Resources}
\begin{itemize}[leftmargin=1.4em,itemsep=2pt]
  \item Ludwig, J., Mullainathan, S., \& Rambachan, A. (2026). \emph{Large Language Models: An Applied Econometric Framework}. Annual Review of Economics, 18, 283--316. \url{https://www.nber.org/papers/w33344} \, / \, \url{https://arxiv.org/abs/2412.07031}
  \item Cunningham, S. \emph{Explainer of Ludwig, Mullainathan and Rambachan's 2026 Econometrics of LLM Paper}. Scott's Mixtape Substack. \url{https://causalinf.substack.com/p/explainer-of-ludwig-mullainathan}
  \item Cunningham, S. \emph{Causal Inference: The Mixtape}. \url{https://mixtape.scunning.com/}
  \item Anthropic. \emph{Claude Code documentation}. \url{https://code.claude.com/docs}
\end{itemize}

\end{document}
